\documentclass[11pt,a4paper]{article}

\usepackage[margin=2.5cm]{geometry}
\usepackage{amsmath,amssymb}
\usepackage{enumitem}
\usepackage[most]{tcolorbox}

% ---------- Notation ----------
\newcommand{\vb}[1]{\mathbf{#1}}
\newcommand{\vect}[1]{
    \begin{pmatrix}#1
\end{pmatrix}}
\newcommand{\PP}{\mathbb{P}}
\newcommand{\RR}{\mathbb{R}}

% ---------- Consistent layout ----------
% Every exercise:  \exercise{n}
% Every sub-question: \part{label}{short description}
% Every final result: \begin{answer} ... \end{answer}
\newcommand{\exercise}[1]{
\section*{Exercise #1}}
\newcommand{\task}[2]{
\subsection*{#1\quad\normalfont\textit{#2}}}
\newtcolorbox{answer}{colback=gray!8,colframe=gray!60,boxrule=0.5pt,
    arc=2pt,left=6pt,right=6pt,top=4pt,bottom=4pt,
fontupper=\normalsize,before upper={\textbf{Answer.}\ }}

\setlength{\parindent}{0pt}
\setlength{\parskip}{6pt}

\title{Assignment 1  Theoretical Exercises}
\author{Computer Vision, 2026 \\ Arvid Kolstad}
\date{}

\begin{document}
\maketitle

% =====================================================================
\exercise{1}

Given the homogeneous points
\[
    \vb{x}_1 = \vect{-4\\3\\1},\qquad
    \vb{x}_2 = \vect{4\\-5\\-2},\qquad
    \vb{x}_3 = \vect{-6\lambda \\12\lambda \\3\lambda}.
\]

\task{(1)}{Cartesian coordinates}

Divide each point by its last coordinate:
\[
    \vb{x}_1 \;\mapsto\; \left(\tfrac{-4}{1},\tfrac{3}{1}\right),\qquad
    \vb{x}_2 \;\mapsto\; \left(\tfrac{4}{-2},\tfrac{-5}{-2}\right),\qquad
    \vb{x}_3 \;\mapsto\;
    \left(\tfrac{-6\lambda}{3\lambda},\tfrac{12\lambda}{3\lambda}\right).
\]
\begin{answer}
    $\vb{x}_1 \text{ maps to } (-4,\,3)$, \quad $\vb{x}_2 \text{ maps to }
    (-2,\,\tfrac{5}{2})$, \quad $\vb{x}_3  \text{ maps to }(-2,\,4)$.
\end{answer}

\task{(2)}{Interpretation of $\vb{x}_4$}

\[
    \vb{x}_4 = \vect{5\\-4\\0}.
\]
Replace the last coordinate by a small number $\varepsilon$. The
corresponding 2D point is
\[
    \left(\frac{5}{\varepsilon},\,\frac{-4}{\varepsilon}\right).
\]
As $\varepsilon \to 0$ this point goes to infinity in the
direction $(5,-4)$.
\begin{answer}
    $\vb{x}_4$ is a point at infinity (a vanishing point) in the
    direction $(5,-4)$.
\end{answer}

\task{(3)}{Are $\vb{x}_4$ and $\vb{x}_5$ the same point?}
\[
    \vb{x}_5 = \vect{5\\-3\\0},\qquad
\]
By the argument in (2), $\vb{x}_4$ goes to infinity in the
direction $(5,-4)$,
while $\vb{x}_5$ goes to infinity in the direction $(5,-3)$. These
directions are not the same.
\begin{answer}
    No, $\vb{x}_4 \neq \vb{x}_5$, since they are points at infinity
    in different directions.
\end{answer}

% =====================================================================
\exercise{2}

\task{(4)}{Intersection of $\vb{l}_1$ and $\vb{l}_2$}

\[
    \vb{l}_1 = \vect{1\\3\\1},\qquad \vb{l}_2 = \vect{2\\-1\\1}.
\]
A point $\vb{x}=(x,y,z)^T$ lies on a line $\vb{l}$ if $\vb{l}^T\vb{x}=0$.
The intersection must lie on both lines:
\[
    \begin{cases}
        \vb{l}_1^T\vb{x} = 0\\
        \vb{l}_2^T\vb{x} = 0
    \end{cases}
    \;\Longrightarrow\;
    \begin{cases}
        x + 3y + z = 0\\
        2x - y + z = 0
    \end{cases}
\]
Set $z=1$. Subtracting the equations gives $-x + 4y = 0$, i.e.\ $x = 4y$.
Inserting into the second equation gives $7y = -1$, so
\[
    y = -\tfrac{1}{7},\qquad x = -\tfrac{4}{7}.
\]
\begin{answer}
    $\vb{x} = \left(-\tfrac{4}{7},\,-\tfrac{1}{7},\,1\right)^T$,
    which in $\RR^2$ is the point
    $\left(-\tfrac{4}{7},\,-\tfrac{1}{7}\right)$.
\end{answer}

\task{(5)}{Intersection of $\vb{l}_3$ and $\vb{l}_4$}

\[
    \vb{l}_3 = \vect{2\\0\\2},\qquad \vb{l}_4 = \vect{-3\\0\\3}.
\]
From the line equation:
\[
    \begin{cases}
        2x + 2z = 0\\
        -3x + 3z = 0
    \end{cases}
    \;\Longrightarrow\;
    \begin{cases}
        x = -z\\
        x = z
    \end{cases}
    \;\Longrightarrow\; x = z = 0,
\]
while $y$ is free.
\begin{answer}
    $\vb{x} = (0,\,1,\,0)^T$. This is a vanishing point: the
    intersection lies at infinity in the $y$-direction.
\end{answer}

\task{(6)}{Line through $x_1$ and $x_2$}

\[
    x_1 = (1,3) \;\Rightarrow\; \vb{x}_1 = \vect{1\\3\\1},\qquad
    x_2 = (2,-1) \;\Rightarrow\; \vb{x}_2 = \vect{2\\-1\\1}.
\]
The line must satisfy $\vb{l}^T\vb{x}_1 = \vb{l}^T\vb{x}_2 = 0$. This
is the same system as in (4),
so we have already calculated it to be $\vb{l} =
\left(-\tfrac{4}{7},-\tfrac{1}{7},1\right)^T$. Check:
\[
    -\tfrac{4}{7} - \tfrac{3}{7} + 1 = 0,\qquad
    -\tfrac{8}{7} + \tfrac{1}{7} + 1 = 0.
\]
This is because in $\PP^2$ the equation is a plane through the origin
that contains the points
$\vb{x}_1$ and $\vb{x}_2$. When collapsing $\PP^2$ to $\RR^2$, the
plane becomes the line between $x_1$ and $x_2$.
\begin{answer}
    $\vb{l} = \left(-\tfrac{4}{7},\,-\tfrac{1}{7},\,1\right)^T$.
\end{answer}

% =====================================================================
\exercise{3}

\task{(7)}{Why the null space gives the intersection}

% NOTE: written as a sum in the original notes; see the comments
% in the chat.
\[
    M\vb{x} =
    [\vb{l}_2^T\vb{x} ,\vb{l}_1^T\vb{x}]^T.
\]
This is just the line equation for the two lines. Since we know that
these must hold for the lines then the vector that gets created must
be the zero vector.
\begin{answer}
    A vector in the null space of $M$ satisfies the line equation of
    both lines, so it is their intersection point.
\end{answer}

\task{(8)}{Are there other solutions?}

Of course $\vb{x}=\vb{0}$ is also solution to this equation, but this
does not have a representation in $\RR^2$ and can therefore be
disregarded for our case.

\begin{answer}
    There are no other points in $\PP^2$ that solve the equation,
    except the zero-vector.
\end{answer}

% =====================================================================
\exercise{4}

\[
    H = \vect{1&1&1\\-1&1&1\\0&0&1},\qquad
    \vb{x}_1 = \vect{1\\0\\1},\qquad
    \vb{x}_2 = \vect{0\\1\\1}.
\]

\task{(a)}{Transformed points}

\[
    \vb{y}_1 = H\vb{x}_1 = \vect{2\\0\\1},\qquad
    \vb{y}_2 = H\vb{x}_2 = \vect{2\\2\\1}.
\]
\begin{answer}
    $\vb{y}_1 = (2,0,1)^T$, \quad $\vb{y}_2 = (2,2,1)^T$.
\end{answer}

\task{(b)}{Lines through the points}

Line $\vb{l}_1 = (a,b,c)^T$ through $\vb{x}_1$ and $\vb{x}_2$:
\[
    \begin{cases}
        \vb{l}_1^T\vb{x}_1 = 0\\
        \vb{l}_1^T\vb{x}_2 = 0
    \end{cases}
    \;\Longrightarrow\;
    \begin{cases}
        a + c = 0\\
        b + c = 0
    \end{cases}
    \;\overset{c=1}{\Longrightarrow}\; a = b = -1.
\]
Line $\vb{l}_2 = (a,b,c)^T$ through $\vb{y}_1$ and $\vb{y}_2$:
\[
    \begin{cases}
        \vb{l}_2^T\vb{y}_1 = 0\\
        \vb{l}_2^T\vb{y}_2 = 0
    \end{cases}
    \;\Longrightarrow\;
    \begin{cases}
        2a + c = 0\\
        2a + 2b + c = 0
    \end{cases}
    \;\overset{c=1}{\Longrightarrow}\; a = -\tfrac{1}{2},\; b = 0.
\]
\begin{answer}
    $\vb{l}_1 = (-1,\,-1,\,1)^T$, \quad $\vb{l}_2 =
    \left(-\tfrac{1}{2},\,0,\,1\right)^T$.
\end{answer}

\task{(c)}{Verify $\vb{l}_2 \sim (H^{-1})^T\vb{l}_1$}

Use $H^{-1} = \frac{1}{\det H}\operatorname{adj}(H)$ with
\[
    \det H = 2,\qquad
    \operatorname{adj}(H) = \vect{1&-1&0\\1&1&-2\\0&0&2}
    \;\Longrightarrow\;
    H^{-1} = \frac{1}{2}\vect{1&-1&0\\1&1&-2\\0&0&2},\quad
    (H^{-1})^T = \frac{1}{2}\vect{1&1&0\\-1&1&0\\0&-2&2}.
\]
Then
\[
    (H^{-1})^T\vb{l}_1
    = \frac{1}{2}\vect{1&1&0\\-1&1&0\\0&-2&2}\vect{-1\\-1\\1}
    = \frac{1}{2}\vect{-2\\0\\4}
    = \vect{-1\\0\\2}
    \sim \vect{-\tfrac{1}{2}\\0\\1}.
\]
\begin{answer}
    $(H^{-1})^T\vb{l}_1 = (-1,0,2)^T \sim \vb{l}_2$, so the two lines
    are the same up to scale.
\end{answer}

% =====================================================================
\exercise{5}

\task{}{Lines transform with $(H^{-1})^T$}

Assume $\vb{x}$ lies on $\vb{l}_1$ and let $\vb{y} \sim H\vb{x}$. Then
\[
    \vb{l}_1^T\vb{x} = 0
    \;\Longrightarrow\;
    \vb{l}_1^T H^{-1}\underbrace{H\vb{x}}_{\vb{y}} = 0
    \;\Longrightarrow\;
    \vb{l}_1^T H^{-1}\vb{y} = \left((H^{-1})^T\vb{l}_1\right)^T\vb{y} = 0.
\]
Therefore, if we set $\vb{l}_2 \sim (H^{-1})^T\vb{l}_1$, the
transformed point $\vb{y} \sim H\vb{x}$ belongs to $\vb{l}_2$.
We can always create the inverse of $H$ because a projective
transformation is always reversible.
\begin{answer}
    $\vb{l}_2 \sim (H^{-1})^T\vb{l}_1$.
\end{answer}

% =====================================================================
\exercise{6}

$H_1 = \vect{1&1&\sqrt{2}\\-1&1&1\\0&0&\sqrt{2}}, H_2 =
\vect{2&2&1\\-1&1&0\\1&0&2},H_3 = \vect{1&-2&1\\0&1&0\\1&0&1},H_4 =
\vect{\frac{1}{\sqrt{2}}&0&1\\1&\frac{1}{\sqrt{2}}&0\\0&0&1},$

\task{(a)}{Which are projective transformations?}

\begin{answer}
    $H_1$, $H_2$ and $H_4$ are projective transformations.
\end{answer}

\task{(b)}{Which are affine transformations?}

\begin{answer}
    $H_1$ and $H_4$ are affine.
\end{answer}

\task{(c)}{Which are similarity transformations?}

rotation matrix.
\begin{answer}
    $H_1$.
\end{answer}

\task{(d)}{Which are Euclidean transformations?}
\begin{answer}
    $H_1$.
\end{answer}

\task{(e)}{Which of the transformations preserve lengths between points?}
\begin{answer}
    Only $H_1$ (Euclidean).
\end{answer}

\task{(f)}{Which map lines to lines?}
\begin{answer}
    $H_1$, $H_2$ and $H_4$
\end{answer}

\task{(g)}{Which map parallel lines to parallel lines?}
\begin{answer}
    $H_1$ and $H_4$.
\end{answer}

% =====================================================================
\exercise{7}

The camera equation is $\lambda\vb{x} = P\vb{X}$, with
\[
    P = \vect{1&0&0&0\\0&1&0&0\\0&0&1&-1},\qquad
    \vb{X}_1 = \vect{1\\-1\\2\\1},\quad
    \vb{X}_2 = \vect{-2\\2\\1\\1},\quad
    \vb{X}_3 = \vect{1\\2\\3\\1}.
\]

\task{(a)}{Projections}

\[
    P\vb{X}_1 = \vect{1\\-1\\1},\qquad
    P\vb{X}_2 = \vect{-2\\2\\0},\qquad
    P\vb{X}_3 = \vect{1\\2\\2}.
\]
The last coordinate of $P\vb{X}_2$ is $0$, so its projection is a
vanishing point.
\begin{answer}
    $\vb{x}_1 \sim (1,-1,1)^T$, \quad $\vb{x}_2 \sim (-2,2,0)^T$
    (vanishing point), \quad $\vb{x}_3 \sim (1,2,2)^T$.
\end{answer}

\task{(b)}{Camera center}

The camera center is the point that projects to $\vb{0}$:
\[
    \vect{0\\0\\0} = \vect{X_1\\X_2\\X_3-1}
    \;\Longrightarrow\; X_1 = 0,\; X_2 = 0,\; X_3 = 1.
\]
\begin{answer}
    The camera center is $C = (0,\,0,\,1)$.
\end{answer}

\task{(c)}{Principal axis}

The points on the principal axis project to the image center $(0,0,1)^T$:
\[
    \vect{0\\0\\1} = \vect{X_1\\X_2\\X_3-1}
    \;\Longrightarrow\; X_1 = 0,\; X_2 = 0,\; X_3 = 2.
\]
So the axis goes through $C=(0,0,1)$ and $(0,0,2)$.
\begin{answer}
    The principal axis is the line through $(0,0,1)$ with direction
    $(0,\,0,\,1)$.
\end{answer}

% =====================================================================
\exercise{8}

Let $\vb{x} \in \PP^2$ and $P_1 = [I\;\vb{0}]$.

\task{(a)}{Points that project to $\vb{x}$}

\[
    \vb{x} \sim P_1\vb{U} = [I\;\vb{0}]\,\vb{U}.
\]
Write $\vb{U} = (x_1',\,x_2',\,x_3',\,s')^T$. Then
\[
    P_1\vb{U} = \vect{x_1'\\x_2'\\x_3'} = \vect{x_1\\x_2\\x_3}
    \;\Longrightarrow\;
    \vb{U} \sim \vect{\vb{x}\\ s}.
\]
The $\vb{U}$ is a collection of points that creates a line in 3 dimensions.
It's not possible to know the value of $s$ from the information of
$P_1$ because the part that multiplies with $s$ is zero in this matrix.
\begin{answer}
    $\vb{U}(s) = (\vb{x}^T,\, s)^T$ projects to $\vb{x}$ for
    every $s \in \RR$, and this is a collection of points in 3
    dimensions that creates a line. It's not possible to know the
    value of $s$ by the matrix $P_1$.
\end{answer}

\task{(b)}{The $s$ that puts $\vb{U}(s)$ in the plane}

Assume $\vb{U}$ belongs to the plane $\Pi = (\pi^T,\,1)^T$:
\[
    \Pi^T\vb{U}(s) = 0
    \;\Longrightarrow\;
    \pi^T\vb{x} + s = 0
    \;\Longrightarrow\;
    s = -\pi^T\vb{x}.
\]
\begin{answer}
    $s = -\pi^T\vb{x}$, \quad i.e.\ \ $\vb{U}(s) =
    \vect{\vb{x}\\ -\pi^T\vb{x}}$.
\end{answer}

\task{(c)}{The homography $H = R - \vb{t}\pi^T$}

Given $\vb{x} \sim P_1\vb{U}$, $\vb{y} \sim P_2\vb{U}$ with $P_2 =
[R\;\vb{t}]$, and $\Pi^T\vb{U} = 0$:
\[
    \vb{y} \sim P_2\vb{U}(s)
    = [R\;\vb{t}]\vect{\vb{x}\\ -\pi^T\vb{x}}
    = R\vb{x} - \vb{t}\pi^T\vb{x}
    = \left(R - \vb{t}\pi^T\right)\vb{x}.
\]
\begin{answer}
    $\vb{y} \sim H\vb{x}$ with $H = R - \vb{t}\pi^T$, so $H$ maps
    $\vb{x}$ to $\vb{y}$. \hfill$\square$
\end{answer}

\end{document}
